\documentclass[12pt]{extarticle}
\usepackage{array}
\usepackage[T1]{fontenc}
\usepackage[Russian]{babel}
\usepackage{booktabs}
\usepackage{graphicx}
\usepackage{spreadtab}
\usepackage{numprint}
\usepackage[margin=3cm]{geometry}
\usepackage{float}
\usepackage{xfp}
\usepackage{amsmath}   
\begin{document}
\author{Осипов Александр Б06-604}
\title{Лабораторная работа 1.1.6. Изучение электронного осцилографа}
\maketitle
\section{Цель}

Ознакомиться с устройством и органами управления электронного или цифрового осциллографа.
Научиться измерять амплитуды и частоты произвольных сигналов, изучить основные характеристиик осциллографа.

\section{Оборудование}
Осцилограф, генераторы электрических сигналов, соединительные кабели.

\section {Теоретическая часть}

\section{Ход работы}

\begin{spreadtab}{{tabular}{c|c|c|c|c|c|c}}
@ $\nu$ (ЗГ) & @ T, дел. & @ TIME/DIV, мс & @T, с & @ $\nu$, кГц & @ $\delta\nu$, Гц & @ $\nu_{\text{ГЗ}} - \nu$ \\
1.0283 & 24 & 0.2 & round(b2/5*c2, 2) & round(1/d2, 2) & round((0.02/d2) *e2, 2) & e2 - a2 \\
18.6*10^(-3) & 26 & 10 & round(b3/5*c3, 2) & round(1/d3, 2) & round((0.02/d3) *e3, 2) & e3 - a3 \\
5.6744 & 17.5 & 50*10^(-3) & round(b4/5*c4, 2) & round(1/d4, 2) & round((0.02/d4) *e4, 2) & e4 - a4 \\
184 * 10^(-3) & 27.5 & 1 & round(b5/5*c5, 2) & round(1/d5, 2) & round((0.02/d5) *e5, 2) & e5 - a5 \\
56.082 & 18 & 5*10^(-3) & round(b6/5*c6, 2) & round(1/d6, 2) & round((0.02/d6) *e6, 2) & e6 - a6 \\
548.74 & 18 & 0.5*10^(-3) & round(b7/5*c7, 2) & round(1/d7, 2) & round((0.02/d7) *e7, 2) & e7 - a7 \\
\end{spreadtab}

\end{document}